\documentclass[11pt]{article}
\usepackage[margin=1in]{geometry}
\usepackage[T1]{fontenc}
\usepackage[utf8]{inputenc}
\usepackage{lmodern}
\usepackage[varqu,scaled=0.95]{zi4}
\usepackage{amsmath,amssymb,amsthm,mathtools}
\usepackage[dvipsnames]{xcolor}
\usepackage{fvextra}
\usepackage{booktabs,enumitem,microtype,needspace}
\usepackage[colorlinks,linkcolor=MidnightBlue,citecolor=MidnightBlue,urlcolor=MidnightBlue]{hyperref}
\usepackage[capitalise,nameinlink]{cleveref}
\newtheorem{theorem}{Theorem}[section]
\newtheorem{lemma}[theorem]{Lemma}
\theoremstyle{definition}
\newtheorem{definition}[theorem]{Definition}
\theoremstyle{remark}
\newtheorem{remark}[theorem]{Remark}
\newcommand{\leanname}[1]{\texttt{\detokenize{#1}}}
\newenvironment{leanblock}[1]
  {\par\smallskip\Needspace{4\baselineskip}\noindent{\footnotesize\textsf{#1}}\par\nopagebreak\vspace{2pt}}
  {\par\medskip}
\fvset{breaklines=true,breakanywhere=true,fontsize=\footnotesize,frame=leftline,
  framerule=0.8pt,rulecolor=\color{gray!60},xleftmargin=4pt,framesep=6pt}
\usepackage{longtable}
\newcommand{\unverified}{\quad\textup{\small[not machine-verified]}}

\newcommand{\Gc}{G_c(\mathfrak{m})}

\title{Vanishing of the unipotent generators at zero in the presented group $G(\mathfrak{m})$}
\date{}
\begin{document}
\maketitle
\begin{abstract}
We prove Lemma 3.1 (book label "gsimple zero generators trivial"). It states that in the group $G(\mathfrak{m})$ given by generators and relations, each unipotent generator $X_{-1}(0)$, $Y_{-1}(0)$, $X_{\ell,jk}(0)$, $Y_{\ell,jk}(0)$ is the identity. In the formalization, the coefficient function $c$ of the modular function $J$ is replaced by an arbitrary function $c\colon\mathbb{N}\to\mathbb{N}$, and the result holds for every such $c$. The proof has two steps. First, a general lemma: a map from $(\mathbb{C},+)$ to a group that turns sums into products sends $0$ to $1$. Second, four lemmas showing that each family of unipotent generators has this property, by the additivity relations (Re:1a), (Re:1b), (Im:1a) and (Im:1b). The main theorem and all lemmas are machine-verified.
\end{abstract}

\begin{center}\small
\begin{tabular}{@{}ll@{}}
\toprule
Source & Carbone et al., G-Simple (draft), GSimple\_Lemma\_3\_1 \\
Status & proof machine-verified (Lean 4, Mathlib) \\
Lemmas & 5 \\
Text & written by claude-opus-5-5, checked against the Lean source \\
Run & \texttt{srv\_20261004\_123233\_GSimple\_Lemma\_3\_1} \\
\bottomrule
\end{tabular}
\end{center}

\tableofcontents

\section{Setting}

Throughout, $\mathbb{N}=\{0,1,2,\dots\}$, $\mathbb{C}^\times$ is the group of units of $\mathbb{C}$, and $c\colon \mathbb{N}\to\mathbb{N}$ is an arbitrary function. The function $c$ plays the role of the coefficient function of $J$; no property of $c$ is assumed.

\begin{definition}[Index set]
Let
\[
E_c=\{(\ell,j,k)\in\mathbb{N}^3 \mid 1\le k\le c(j),\ \ell<j\}.
\]
Since $\ell<j$, every $(\ell,j,k)\in E_c$ has $j\ge 1$.

For $e=(\ell,j,k)\in E_c$, put $e^{\ast}=(j-1-\ell,\,j,\,k)$. Then $e^{\ast}\in E_c$, because $j-1-\ell<j$.
\end{definition}

\begin{definition}[Structure constants]
For $\ell,j\in\mathbb{N}$, let
\[
c_{\ell j}=(-1)^{\ell+1}\binom{j-1}{\ell}(\ell+1)(j-\ell)\in\mathbb{C}.
\]
\end{definition}

For $\ell<j$ all differences here are ordinary differences of natural numbers, and $c_{\ell j}\neq 0$. Divisions by $c_{\ell j}$ below are therefore genuine divisions in $\mathbb{C}$.

\begin{definition}[Generators]
Let $\mathcal{X}_c$ be the set of formal symbols
\[
H_1(s),\ H_2(s)\ (s\in\mathbb{C}^\times),\qquad X_{-1}(u),\ Y_{-1}(u)\ (u\in\mathbb{C}),\qquad X_e(u),\ Y_e(u)\ (e\in E_c,\ u\in\mathbb{C}).
\]
For $e=(\ell,j,k)$ we also write $X_{\ell,jk}(u)=X_e(u)$ and $Y_{\ell,jk}(u)=Y_e(u)$. Let $F(\mathcal{X}_c)$ be the free group on $\mathcal{X}_c$.

In $F(\mathcal{X}_c)$ define, for $s\in\mathbb{C}^\times$ and $e=(\ell,j,k)\in E_c$:
\[
\widetilde w_{-1}(s)=X_{-1}(s)\,Y_{-1}(-s^{-1})\,X_{-1}(s),\qquad \widetilde w_{-1}=\widetilde w_{-1}(1),
\]
\[
\widetilde w_e(s)=X_e(s)\,Y_e\!\left(\frac{-s^{-1}}{c_{\ell j}}\right)X_e(s),\qquad \widetilde w_e=\widetilde w_e(1).
\]
The commutator is $(a,b)=aba^{-1}b^{-1}$.
\end{definition}

\begin{definition}[Relators]
Let $\mathcal{R}_c\subseteq F(\mathcal{X}_c)$ consist of the following elements. In each list we write a relation ``$a=b$'', which stands for the relator $ab^{-1}$. A relation ``$(a,b)=1$'' stands for the relator $(a,b)$. The variables range over $s,t\in\mathbb{C}^\times$, $u,v\in\mathbb{C}$, and $e=(\ell,j,k)$, $f=(m,p,q)\in E_c$.
\begin{itemize}
\item (H:1a) $H_1(s)H_1(t)=H_1(st)$.
\item (H:1b) $H_2(s)H_2(t)=H_2(st)$.
\item (H:2) $H_1(s)H_2(t)=H_2(t)H_1(s)$.
\item (Re:1a) $X_{-1}(u)X_{-1}(v)=X_{-1}(u+v)$.
\item (Re:1b) $Y_{-1}(u)Y_{-1}(v)=Y_{-1}(u+v)$.
\item (Re:2) $Y_{-1}(-t)X_{-1}(s)Y_{-1}(t)=X_{-1}(-t^{-1})\,Y_{-1}(-t^2s)\,X_{-1}(t^{-1})$.
\item (Re:3) $\widetilde w_{-1}(s)\widetilde w_{-1}=H_1(-s)H_2(-s^{-1})$.
\item (Re:4a) $\widetilde w_{-1}X_{-1}(u)\widetilde w_{-1}^{-1}=Y_{-1}(-u)$.
\item (Re:4b) $\widetilde w_{-1}Y_{-1}(u)\widetilde w_{-1}^{-1}=X_{-1}(-u)$.
\item (Re:5a) $\widetilde w_{-1}H_1(s)\widetilde w_{-1}^{-1}=H_2(s)$.
\item (Re:5b) $\widetilde w_{-1}H_2(s)\widetilde w_{-1}^{-1}=H_1(s)$.
\item (Re:6a) $H_1(s)X_{-1}(u)H_1(s)^{-1}=X_{-1}(su)$.
\item (Re:6b) $H_2(s)X_{-1}(u)H_2(s)^{-1}=X_{-1}(s^{-1}u)$.
\item (Re:6c) $H_1(s)Y_{-1}(u)H_1(s)^{-1}=Y_{-1}(s^{-1}u)$.
\item (Re:6d) $H_2(s)Y_{-1}(u)H_2(s)^{-1}=Y_{-1}(su)$.
\item (Im:1a) $X_e(u)X_e(v)=X_e(u+v)$.
\item (Im:1b) $Y_e(u)Y_e(v)=Y_e(u+v)$.
\item (Im:2) $Y_e(-t)X_e(s)Y_e(t)=X_e\!\left(\frac{-t^{-1}}{c_{\ell j}}\right)Y_e(-c_{\ell j}t^2s)\,X_e\!\left(\frac{t^{-1}}{c_{\ell j}}\right)$.
\item (Im:3) $\widetilde w_e\big(s^{(\ell+1)(j-\ell)}\big)\widetilde w_e=H_1\big((-s)^{j-\ell}\big)H_2\big((-s)^{\ell+1}\big)$.
\item (Im:4a) $\widetilde w_eX_e(u)\widetilde w_e^{-1}=Y_e(-u/c_{\ell j})$.
\item (Im:4b) $\widetilde w_eY_e(u)\widetilde w_e^{-1}=X_e(-c_{\ell j}u)$.
\item (Im:5a) $\widetilde w_eH_1(s^{j-\ell})\widetilde w_e^{-1}=H_2(s^{-(\ell+1)})$.
\item (Im:5b) $\widetilde w_eH_2(s^{\ell+1})\widetilde w_e^{-1}=H_1(s^{-(j-\ell)})$.
\item (Im:6a) $H_1(s)X_e(u)H_1(s)^{-1}=X_e(s^{\ell+1}u)$.
\item (Im:6b) $H_2(s)X_e(u)H_2(s)^{-1}=X_e(s^{j-\ell}u)$.
\item (Im:6c) $H_1(s)Y_e(u)H_1(s)^{-1}=Y_e(s^{-(\ell+1)}u)$.
\item (Im:6d) $H_2(s)Y_e(u)H_2(s)^{-1}=Y_e(s^{-(j-\ell)}u)$.
\item (U:1a) $(X_{-1}(u),X_e(v))=1$ if $\ell+1=j$.
\item (U:1b) $(Y_{-1}(u),Y_e(v))=1$ if $\ell+1=j$.
\item (U:1c) $(Y_{-1}(u),X_e(v))=1$ if $\ell=0$.
\item (U:1d) $(X_{-1}(u),Y_e(v))=1$ if $\ell=0$.
\item (U:2) $(X_e(u),Y_f(v))=1$ if $j\ne p$, or $k\ne q$, or $|\ell-m|>1$.
\item (U:3a) $(X_e(u),Y_f(v))=X_{-1}(uv)$ if $e=(1,2,k)$ and $f=(0,2,k)$ for the same $k$.
\item (U:3b) $(X_e(u),Y_f(v))=Y_{-1}(uv)$ if $e=(0,2,k)$ and $f=(1,2,k)$ for the same $k$.
\item (U:4a) $\widetilde w_{-1}X_e(u)\widetilde w_{-1}^{-1}=X_{e^\ast}\big((-1)^{j-\ell-1}u\big)$.
\item (U:4b) $\widetilde w_{-1}Y_e(u)\widetilde w_{-1}^{-1}=Y_{e^\ast}\big((-1)^{j-\ell-1}u\big)$.
\end{itemize}
\end{definition}

\begin{definition}[The group]
Let $\Gc=F(\mathcal{X}_c)/N_c$, where $N_c$ is the normal closure of $\mathcal{R}_c$. Let $\pi\colon F(\mathcal{X}_c)\to\Gc$ be the quotient homomorphism.

By construction, $\pi(r)=1$ for every $r\in\mathcal{R}_c$. Consequently, if $ab^{-1}\in\mathcal{R}_c$, then $\pi(a)=\pi(b)$.
\end{definition}

\begin{theorem}\label{thm:main}
Let $c\colon\mathbb{N}\to\mathbb{N}$ be any function. Then in $\Gc$:
\begin{enumerate}
\item[(a)] $\pi(X_{-1}(0))=1$;
\item[(b)] $\pi(Y_{-1}(0))=1$;
\item[(c)] $\pi(X_e(0))=1$ for every $e\in E_c$;
\item[(d)] $\pi(Y_e(0))=1$ for every $e\in E_c$.
\end{enumerate}
\end{theorem}

This is Lemma 3.1 of the source for an arbitrary coefficient function $c$. Taking $c(j)$ to be the coefficients of $J$ recovers the source statement. The hypothesis $c(2)\ge 1$, mentioned in the source's formalization note, is not needed here.

\section{Lemmas}

\subsection{A general fact about additive families}

\begin{lemma}[Additive families vanish at zero]\label{lem:add_zero_group_lemma}
Let $G$ be a group and $f\colon\mathbb{C}\to G$ a map with $f(u+v)=f(u)f(v)$ for all $u,v\in\mathbb{C}$. Then $f(0)=1$.
\end{lemma}

\begin{proof}
Taking $u=v=0$ gives $f(0)=f(0+0)=f(0)f(0)$. Hence
\[
f(0)=f(0)\big(f(0)f(0)^{-1}\big)=\big(f(0)f(0)\big)f(0)^{-1}=f(0)f(0)^{-1}=1.\qedhere
\]
\end{proof}

\subsection{Additivity of the unipotent families in $\Gc$}

\begin{lemma}[Additivity of $X_{-1}$]\label{lem:mk_xm_add}
For every $c\colon\mathbb{N}\to\mathbb{N}$ and all $u,v\in\mathbb{C}$, we have $\pi(X_{-1}(u+v))=\pi(X_{-1}(u))\,\pi(X_{-1}(v))$ in $\Gc$.
\end{lemma}

\begin{proof}
By (Re:1a), the element $X_{-1}(u)X_{-1}(v)X_{-1}(u+v)^{-1}$ lies in $\mathcal{R}_c$, so its image under $\pi$ is $1$. Since $\pi$ is a homomorphism, this reads
\[
\pi(X_{-1}(u))\,\pi(X_{-1}(v))\,\pi(X_{-1}(u+v))^{-1}=1.
\]
Multiplying on the right by $\pi(X_{-1}(u+v))$ gives the claim.
\end{proof}

\begin{lemma}[Additivity of $Y_{-1}$]\label{lem:mk_ym_add}
For every $c\colon\mathbb{N}\to\mathbb{N}$ and all $u,v\in\mathbb{C}$, we have $\pi(Y_{-1}(u+v))=\pi(Y_{-1}(u))\,\pi(Y_{-1}(v))$ in $\Gc$.
\end{lemma}

\begin{proof}
By (Re:1b), the element $Y_{-1}(u)Y_{-1}(v)\,Y_{-1}(u+v)^{-1}$ lies in $\mathcal{R}_c$, so its image under $\pi$ is $1$. Hence
\[
\pi\big(Y_{-1}(u)Y_{-1}(v)\big)\,\pi(Y_{-1}(u+v))^{-1}=1.
\]
Multiplying on the right by $\pi(Y_{-1}(u+v))$ gives $\pi(Y_{-1}(u)Y_{-1}(v))=\pi(Y_{-1}(u+v))$. Since $\pi$ is a homomorphism, the left side equals $\pi(Y_{-1}(u))\pi(Y_{-1}(v))$.
\end{proof}

\begin{lemma}[Additivity of $X_e$]\label{lem:mk_xx_add}
For every $c\colon\mathbb{N}\to\mathbb{N}$, every $e\in E_c$ and all $u,v\in\mathbb{C}$, we have $\pi(X_e(u+v))=\pi(X_e(u))\,\pi(X_e(v))$ in $\Gc$.
\end{lemma}

\begin{proof}
By (Im:1a), the element $\big(X_e(u)X_e(v)\big)X_e(u+v)^{-1}$ lies in $\mathcal{R}_c$, so $\pi(X_e(u)X_e(v))=\pi(X_e(u+v))$. Since $\pi$ is a homomorphism,
\[
\pi(X_e(u+v))=\pi(X_e(u)X_e(v))=\pi(X_e(u))\,\pi(X_e(v)).\qedhere
\]
\end{proof}

\begin{lemma}[Additivity of $Y_e$]\label{lem:mk_yy_add}
For every $c\colon\mathbb{N}\to\mathbb{N}$, every $e\in E_c$ and all $u,v\in\mathbb{C}$, we have $\pi(Y_e(u+v))=\pi(Y_e(u))\,\pi(Y_e(v))$ in $\Gc$.
\end{lemma}

\begin{proof}
By (Im:1b), the element $\big(Y_e(u)Y_e(v)\big)Y_e(u+v)^{-1}$ lies in $\mathcal{R}_c$, so $\pi(Y_e(u)Y_e(v))=\pi(Y_e(u+v))$. Since $\pi$ is a homomorphism,
\[
\pi(Y_e(u+v))=\pi(Y_e(u)Y_e(v))=\pi(Y_e(u))\,\pi(Y_e(v)).\qedhere
\]
\end{proof}

\section{Proof of the theorem}

\begin{proof}[Proof of \cref{thm:main}]
Fix $c$. Each part applies \cref{lem:add_zero_group_lemma} with $G=\Gc$ to a suitable map $f\colon\mathbb{C}\to\Gc$.
\begin{itemize}
\item[(a)] Take $f(u)=\pi(X_{-1}(u))$. This map is additive by \cref{lem:mk_xm_add}, hence $\pi(X_{-1}(0))=1$.
\item[(b)] Take $f(u)=\pi(Y_{-1}(u))$. This map is additive by \cref{lem:mk_ym_add}, hence $\pi(Y_{-1}(0))=1$.
\item[(c)] Fix $e\in E_c$ and take $f(u)=\pi(X_e(u))$. This map is additive by \cref{lem:mk_xx_add}, hence $\pi(X_e(0))=1$.
\item[(d)] Fix $e\in E_c$ and take $f(u)=\pi(Y_e(u))$. This map is additive by \cref{lem:mk_yy_add}, hence $\pi(Y_e(0))=1$.\qedhere
\end{itemize}
\end{proof}

\paragraph{Relation to the source proof.}
The source proof treats only part (a). It sets $u=v=0$ in (Re:1a) to get $X_{-1}(0)^{-1}X_{-1}(0)X_{-1}(0)=1$, and states that (b), (c), (d) follow in the same way from (Re:1b), (Im:1a) and (Im:1b).

The argument here is the same in substance, organised in two steps. First, it extracts from each additivity relation the statement that $u\mapsto\pi(\cdot(u))$ is a homomorphism $(\mathbb{C},+)\to\Gc$. Then it applies, in all four cases, the general fact that such a map sends $0$ to $1$. All four parts are carried out explicitly.

The only other difference is the one noted after the theorem. The coefficient function $c$ is an arbitrary parameter, and no hypothesis on $c$ is used, since only the relations (Re:1a), (Re:1b), (Im:1a) and (Im:1b) enter the proof.

\appendix
\section{Correspondence with the formal proof}
Each lemma of the text and the Lean declaration that states it; \cref{thm:main} is \texttt{gsimple\_zero\_generators\_trivial}.

\begin{longtable}{@{}lll@{}}
\toprule
Lemma & Lean declaration & Status \\
\midrule
\endhead
\cref{lem:add_zero_group_lemma} & \texttt{add\_zero\_group\_lemma} & verified \\
\cref{lem:mk_xm_add} & \texttt{mk\_xm\_add} & verified \\
\cref{lem:mk_ym_add} & \texttt{mk\_ym\_add} & verified \\
\cref{lem:mk_xx_add} & \texttt{mk\_xx\_add} & verified \\
\cref{lem:mk_yy_add} & \texttt{mk\_yy\_add} & verified \\
\bottomrule
\end{longtable}
\end{document}
